% Benchmark paragraph categories matching the paper (Table 1): short (1 line),
% medium (4-5 lines), long (10+ lines), inline math, display math.
\long\def\parShort{The quick brown fox jumps over the lazy dog.}
\long\def\parMedium{Typesetting a paragraph is a local operation: given the same
measure and the same content, the line breaker produces identical breaks whatever
precedes or follows it. Word processors have relied on this for decades, reflowing
only the paragraph under the cursor while a background process reconciles the
global layout of the document at its own pace.}
\long\def\parLong{The line-breaking algorithm considers the paragraph as a whole,
choosing breakpoints that minimise the total demerits accumulated across all of its
lines rather than greedily filling each line in turn. Because the optimisation is
confined to a single paragraph, the result does not depend on any other paragraph
in the document. Page breaking, by contrast, is global: changing the height of one
paragraph can move every subsequent line, float and footnote, shifting page numbers
and cross-references along the way. The insight behind real-time editing is that
only the first operation needs to happen on every keystroke. The second can run in
the background and converge within seconds, exactly as a word processor does. With
microtypographic extensions enabled the engine also expands or shrinks glyphs by a
few per mille and lets punctuation protrude into the margin, which must be captured
faithfully by any renderer that bypasses the usual PDF output. All of this happens
in well under two milliseconds for a paragraph of this length on a modern laptop.}
\long\def\parInlineMath{Let $f(x) = \sum_{n=0}^{\infty} a_n x^n$ converge for $|x| < R$,
where $R = 1/\limsup |a_n|^{1/n}$ is the radius of convergence of the series.}
\long\def\parDisplayMath{\[ \int_{-\infty}^{\infty} e^{-x^2}\,dx = \sqrt{\pi} \]}
